\documentclass[11pt,a4paper]{article}

% ENCODING AND TYPOGRAPHY
\usepackage[T1]{fontenc}
\usepackage[utf8]{inputenc}
\usepackage{libertine}
\usepackage{microtype}

% PAGE LAYOUT
\usepackage[a4paper,
            top=1.5cm,
            bottom=1.5cm,
            left=1.5cm,
            right=1.5cm]{geometry}

% DUMMY TEXT
\usepackage{lipsum}


\begin{document}

\noindent
{\small
\textbf{Author:} HENNINGSEN, Charles Frederick\\
\textbf{Year:} 1846\\
\textbf{Title:} "Finland and Its Literature"\\
\textbf{Source:} Eastern Europe and the Emperor Nicholas, II\\
\textbf{Publisher:} T. C. Newby\\
\textbf{Pages:} 154--210
}

\vspace{1em}

CHAPTER VIII

FINLAND AND ITS LITERATURE

**Swedish feeling in Finland - Attempt of the Russian Cabinet to neutralise it by the resuscitation of a Finnish Nationality and Literature - Famous Epic Poem of the Kalevala - Collection of the Kanletetar.**

T[he] principality of Finland, inhabited by an aboriginal race, extends from within a few hours' sail of Stockholm up to the gates of St. Petersburgh.

Its population received its laws, civilisation, and religion from the Swedes. Its higher classes are all of Swedish origin; but as complete an amalgamation has taken place as between the English population in Wales, or the French amongst the Celts of Britany.

The Swedish language is used in precisely the same manner as the English or the French in those Celtic countries. Finland was, as it is well known, lost to Sweden under the last but one of the dynasty of Vasas, an obtuse and hasty, but wellmeaning sovereign.

It was occupied by Alexander when he had been forced by Napoleon to join with him against the allies with whom he had begun the war, and lost to Sweden by the only monarch who remained faithful, in the continental struggle, to his engagements with Great Britain.

Alexander retained possession of Finland, notwithstanding his private assurance at the time that he had no intention of appropriating it.

The Finns have never ceased to regard with affectionate regret the period of their union with Sweden. In Sweden, the loss of Finland rankles so profoundly in every Swedish heart, that from Gottenburg to the borders of Lapland you cannot name the subject to the humblest peasant in the remotest hut, without seeing that the prospect of its recovery would reconcile him to any sacrifice.

Finland, from the origin of one part of its population and the predilection of the remainder, from the nature of its surface, so chequered by rocks, lakes, defiles, forests, and rivers, and peculiarly unfitted for the action of Muscovite masses, is the most favourable ground the Swedes could choose the theatre of a successful struggle.

Sweden is every day more likely to increase its power by incorporating the sister population of Denmark (with its comparative wealth and excellent sailors) in a common union.

The Danish and Swedish people, long divided by the consequences of a family feud, are identical in race and language; and their interests are the same. The prejudice which blinded them to this fact has recently been rapidly clearing away, and they are awakening to a conviction of the truth. It is to be doubted whether at the present day any political combination can eventually prevent two neighbouring people, bent on fraternising, from uniting. This Scandinavian union of Norway, Sweden, and Denmark, which the author believes first to have noticed publicly in the Revelations of Russia, and which will become imminent on the death of the present king of Denmark, is gradually occupying more attention in the north. The students of the universities go across in increased numbers to visit each other; and the consummation so much to be wished for has become a universal toast.

In Denmark, where all public demonstrations have rather reference to the destiny of that kingdom after the death of its present sovereign than to any violent changes during his lifetime, a growing spirit of liberalism manifests itself. It is only a few weeks since that a journalist being condemned for a severe article against the absolute monarch of Prussia to a fine of near £500, on the seizure-sale of his library to defray the amount, a gentleman stepped forward and bid the amount for the first book put up for sale.

The eventual results of a Scandinavian union, must undoubtedly lead to an attempt at the recovery of Finland by Sweden; but should its accomplishment be long delayed, the feeling of the Swedes upon this subject furnishes Great Britain with the means, in the event of war with Russia, of dealing at slight cost or hazard, a serious and irrecoverable blow at that empire.

Whatever may be the political condition of Sweden, if, in event of war between Russia and Great Britain, the prospect of the recovery of Finland were held out to its people, no government could resist the popular movement, which would impel it to join with England, and invade that principality. A British fleet, and a Swedish army subsidised by England, could undoubtedly effect the conquest of this territory, in spite of every effort the Russian cabinet could make; and having occupied, permanently retain it.

Cronstadt is built on a Finnish island, and though fortified with every care, being wholly artificial, can make no permanent resistance against the newly introduced means of destruction at the disposal of a victorious navy. The only defence of St. Petersburg, is Cronstadt; and its occupation, together with that of Finland, leave the capital itself at the invaders' mercy, and must lead to its abandonment.

War with Great Britain would, besides, at any moment rouse Poland, which is hourly drawing nearer to the renewal of a single-handed struggle with the Tsar; Turkey would at all times be eager to profit by such an eventuality; an English fleet in the Black Sea could play the same part as in the Baltic, and a small military expediton might deprive Russia of her Georgian provinces, from which she is almost cut off by the Caucasian mountaineers, and where it is four times more difficult and expensive for her to send a soldier, than it would be to Great Britain.

The author, who is prepared to enter elsewhere into military details to prove the practicability of the aggressive means to which he makes here allusion, wishes it to be understood that he is no partisan of war, which cannot in his estimation be sanctioned on any grounds of national interest, jealousy, or aggrandisement; which even the pretence of injured national honour will not justify, and which ought never to be undertaken except in self-defence, or for the common-weal of humanity.

There are many with whom this consideration would weigh in their judgment of the conduct to be pursued by Great Britain towards the Russian despot, if it were not for the popular idea that he is as inaccessible to us as the tyrant of Bokhara. It is to assist in removing this strange misconception that the author is led to dwell on this vulnerable point which the mutual relations of Finland and Sweden offer in the panoply of that gigantic Russian despotism, to which might be so appropriately applied the epithet conferred in the whispered murmurs of the crowd, on its imperial representative when passing through Germany, of *Der Menschenfeind*, or the "enemy of humanity."

On the national feeling of the Finns of themselves, whose character the author has sketched elsewhere - a patient, hardy race, distinguished for loyalty and fidelity - naturally depends much of the facility which an invasion of their principality might offer. If they had been as well treated by their new masters as by the Swedes, the course of time would have reconciled them to the change. The author has had occasion to show elsewhere how Alexander and Nicholas swore successively the maintenance of their constitution, and how both violated this oath by never allowing a representative body to assemble. This circumstance recails too another instance of the misrepresentation prevalent in the public press, out of compliment or through fear of the Russian cabinet.

A Monsieur Leouzon le Duc, a modest and timid writer, has recently published a work on Finland, in which he gives a translation of the Kalevala, whether from the Finnish, or from the Swedish translation of Castren he does not mention. In his account of the government of Finland, he states that Alexander promised to mantain the Finnish constitution, and cites his manifest to that effect; he further mentions with apparent approbation, that Nicholas graciously confirmed it at his accession to the throne, and also gives the text of his ukase.

But he utterly omits to say that this promise was never fulfilled - the grievance of which the Finns so bitterly complain; thus naturally leaving his reader to infer that these two sovereigns had respected, instead of virtually abrogating the Finnish constitution, in the face of a pledge so solemn. Alexander otherwise pursued towards that country a system of conciliation; the exemption from duties was guaranteed for a certain number of years, and the vexatious tyranny of the bureaucracy and police restrained within tolerable bounds. The impatience of Nicholas soon began to tighten the rein. Directly after his accession, he abolished the committee for Finnish affairs, and from that moment the country was gradually Russianised, that is to say, that all the corrupt and arbitrary mechanism by which the remainder of the empire is administered, was introduced by slow degrees.

The long peace succeeding to so many years of war, the exemption from duties, and the trade with the capital, consoled the Finns for their political restrictions by an unusual increase of material prosperity. But this has been over-rated; and if we compare the progress of semi-republican Norway, a poorer country, divested of many of the same advantages as Finland, we may infer that this would have been at least the same if still united to Sweden. On the other hand, it must not be forgotten that the exemption from heavy duties, by rendering it a vast free port, was the chief cause of its increasing wealth; but these were only guaranteed for a certain term of years.

This term expired since the author had last occasion to advert to Finland, and as he prognosticated, the first step taken by Nicholas was to deprive it of these advantages. After taking this step, which practically taught his Finnish subjects the value of the representative forms withheld from them, it was discovered that the expense of a sufficient custom-house establishment along a coast of a thousand miles, would exceed the incomings of the customs, and the rigidity of the tariff was again relaxed. This has not been the case with the increasing severity of its administration, and aware of the danger of indisposing this population, the plan has been devised in the imperial cabinet, of resuscitating and playing off the nationality of the Finns against that of the Swedish middle-class, and of neutralizing thereby the Swedish feeling, still so obstinately prevalent in that country.

It is thus, that on the other 'side of the Finnish gulph, the Lettonian peasantry, and their German priesthood and nobility, are skilfully made to keep each other in check. But there is only an apparent similitude in the two cases. In the Baltic provinces, the Teutonic knights reduced the conquered inhabitants to private servitude, and shared them out like cattle; as in parts of Ireland, a fierce animosity lurks in the bosoms of the peasantry against this strange race of oppressors, and may be easily worked upon.

Finland, on the contrary, was equitably governed as a Swedish province, -all the recollections of her civilisation and her glories are identified in her happy association with the mother country. It is therefore as idle to attempt to detach the affections of Finland from Sweden, by reference to a nationality lost in the night of time, or to traditions connected with paganism, as it would be for the French, if in possession of Cornwall, to make its inhabitants forget their separation from England, by encouraging homely songs in their native dialect, or popularizing a poem filled with the mythologic legends of the Druids.

That the reader may be enabled to judge of the nature of Finnish national literature, a brief analysis with extracts from the Kalevala and Kanteletar is herewith given, -the one an ancient epic in two-and-thirty books, the other a collection of the lyrics of the people. The former, which is in many respects a singular and interesting production, may gratify, as a literary curiosity, at the same time that it serves to illustrate the author's argument.

The Kalevala - so called either from Kawa the father of the Finnish gods, the gigantic Saturn of its mythology, or from Kalewala the Finnish Olympus - is a collection of poetic myths, gathered like the songs of Homer, or as those of Ossian were once supposed to have been by Macpherson, from the oral traditions of the people, by the indefatigable Doctor Lœnnroth, and published by him a few years since. The antiquity of the whole of it is very great, as it winds up with the legend of Mariatta the gentle virgin, who swallows a berry, and conceiving, brings into the world a child, which usurps the empire of the god Wainamoinen.

It is impossible in this to mistake the Virgin Mary, and the first tradition of Christianity, dispelling the darkness of heathenism; and there. fore it is obvious that the last portions of the poem can have been composed no later than the twelfth century, when Christianity was introduced into Finland. There are allusions made in other parts of it to the invasions of Russia by the Northmen, which occurred in the ninth century; and some even see cause to retrace the origin of the first portions of it beyond the birth of Christ.

There is, of course, no substantial reason to be given, why these traditions, orally handed down for nine centuries as they contain incontestible proofs of having been, may not in a like manner have commenced and been transmitted from mouth to mouth nine centuries previously.

This poem of the Kalevala, though singularly disjointed, obscure, and confused, may be termed an epic, because all the episodes which it contains refer to the adventures of the god Wainamoinen, of which the thread is never entirely lost from the commencement to the termination. It is, however, an epic almost wholly mythologic, or at least human personages are only casually introduced, and chiefly in the character of sorcerers.

It is compared to the Odyssey by some of the learned Finns, who in their patriotic enthusiasm even give it precedence. Though it is difficult for any one but a Finn to acquiesce in this judgment, the Kalevala must be admitted as an interesting monument, of a distant age and of a numerous people, which probably overspread a large portion of Europe long before it was occupied by its present races, and of which the only notable remnant is now to be found in Finland, though slender branches separated from it, and plunged in barbarism, still exist, scattered over Siberia and the north-east of Europe. This poem is also peculiar, and original, though in some passages its style is biblical, and seems to point to an oriental origin.

It is therefore not without striking beauties, and many more are probably lost to the modern reader, through the unintelligibility of the allegories and allusions with which it is filled, whilst in a translation the harmony of the versification, one of its chiefest charms, is necessarily lost upon a stranger. It is impossible to give any outline of the construction of this poem, without some slight explanation of the Finnish mythology, which is peculiar, though bearing some traces of having been derived either from that of the Greeks, or from a common source. We must however admit it to be less poetical than that of the Greeks, and less gloomily grand than that of the Scandinavians.

It mixes up the weird with the terrible, and though we have no right to condemn the most gross apparent absurdities, which may be full of meaning lost to us, its most solemn superstitions are often chequered by images irresistibly ludicrous.

All the Finnish gods, whenever their deeds or history are recorded, appear so frequently in the character of sorcerers, that the attributes of the godhead seem inextricably confounded with the power of the magician.

They blend in a far greater degree than the gods of the Greek mythology, or of the Scandinavian Edda, the weaknesses of the man with their god-like character; rather wizards than deities, the art with which they use their spells and incantation seems chiefly to enable them to struggle with, and triumph over mortals and inferior spirits.

There is, besides, so much of contradiction and obscurity in all the Finnish myths; there appear to be so many words of which the signification is either changed or lost, that a perfectly distinct idea cannot be formed with certainty, even of the attributes of the principal deities, respecting which, the notions of the singers of the runes themselves differ perhaps as much as those of their modern commentators.

Jumala is the god of clouds and thunder; Wainamoinen of poetry and music. They represent the Jupiter and Apollo, but it is uncertain which of the two was the supreme being.

Kava the giant, the father of the gods and giants, bears some analogy to Saturn. Illmarinen the eternal blacksmith to the Vulcan of the Greeks and Romans. Tuoni is the god of death, the giant Hisii, of evil; Akto is the king of the waters; Tuopio, of the woods; Matha-Teppo, of the roads. The storm is represented by an eagle. Mehilainen is a beneficent bird, small and frail as the humming-bird, but ever bearing on its tiny wings the balm and antidote for sickness, suffering, and the spells of evil.

Besides these, the mighty sorcerers and wizards, there are numerous other gods; and every lake, stream, hill, or valley - in fact, all animate or inanimate things have their good and evil spirits.

On this account everything is personified in their mythic poety. The boat laments upon the shore - the lonely tree, isolated in the clearing, mourns and complains, the road converses with the god, the iron in the furnace has a voice, and in its uses a volition.

But besides these passages, full of originality and beauty, we find the witch of Pohja, whose spells can baffle the gods, sweeping up the dust upon her floor into a brazen pan. The god Wainamoinen cocks his hat gaillardly upon one side, he laves his thumb and purifies his fingers before astonishing all nature with the harmony of his song.

The goddess of the woods draws on her blue stockings, and arrays herself in red ribbons, when irresistibly attracted to listen to the melody of the god.

Illmarinen, another divinity, in the spirit of Hudibras, seriously declares that a man is safer cased in steel, and that he prefers to go by land rather than risk himself upon the water.

It may be asked whether Nicholas or his advisers think these images calculated to efface from the minds of the Finns the recollection of the freedom they enjoyed, and the glories they achieved, when, united to the Swedes, they fought in the front ranks of the armies of the great Gustav Adolph for the civil and religious liberty of thE north.

The Kalevala is divided into two-and-thirty *runes*. Like all Finnish poetry, both ancient and modern, its verse is alliterative instead of being in blank or rhymed; but from the nature of the language, this singular kind of versification is full of harmony; and, to judge from the facility with which it is improvised by the peasantry, must be of easy construction. For instance, the following passage from the first rune of the Kalevala, where the Lapland wizard lies in wait for the god Wainamoinen, and shoots the blue elk on which he is riding:

<div style="font-style: italic">

And he shot his shaft; but it rose too high. The sky was rent, the arches of the air were shaken. He shot a second; but it fell too low: it sunk into the depths of earth-the mother of mendown to *Manala*, whose vaults it made to tremble, &c.

</div>

It would be versified in the Finnish by making most of the words in one line begin with the same letter. As for instance:

<div style="font-style: italic">

Shooting a shaft it soared so high,
That it smote the sky and severed its arch;
Then speeding a second it sunk too low,
And deeply descending down through the dark earth -
Of mortal man the mighty mother -
Made Manalas murmuring vaults to quiver;
Then, with truer aim, his third shaft he shot through 
The blue bounding elk by the hero bestridden.

</div>

The first rune of the Kalevala opens with the birth of the god Wainamoinen, after having been thirty years imprisoned in his mother's womb. He immediately creates himself a courser-a blue and gigantic elk-on which he rides towards the seashore. But there a Lapland wizard lies in ambush for him; and discharging his fatal arrows, kills his strange steed; and the dismounted god having fallen into the sea, wanders for seven years about the waters.

\end{document}